% ============================================================
% Resultate (Handbuch S. 84)
%
% Nur Ergebnisse klar darstellen, hier noch NICHT diskutieren.
% Jede Abbildung/Tabelle braucht eine Legende und wird im Fliesstext
% kommentiert (Handbuch S. 88 f.).
%
% Für dieses Thema z.B.:
% - Messwerte aus dem Brauprozess (Stammwürze/°Plato, pH-Wert,
%   Vergärungsgrad, errechneter Alkoholgehalt)
% - Zeitlicher Verlauf der Gärung (z.B. als Tabelle/Abbildung)
% - Verkostungsergebnisse (falls eine Umfrage/Degustation gemacht wurde)
%
% Tabelle und Diagramm unten werden beim Kompilieren direkt aus
% daten/gaerung.csv erzeugt (pgfplotstable/pgfplots) -- die Messwerte
% stehen NUR in der CSV-Datei, nicht im .tex-Quellcode. Bei neuen
% Messungen einfach die CSV aktualisieren.
%
% Checkliste S. 103: Tabellen müssen auf eine Seite passen, zentriert,
% umrandet und beschriftet (Überschrift oberhalb) sein.
% ============================================================

\chapter{Resultate}
\label{kap:resultate}

% Die Messzeitpunkt-Spalte steht in der CSV im Rohformat
% "YYYY-MM-DD HH:MM" und wird beim Kompilieren auf "DD.MM.YYYY HH:MM"
% umformatiert (Spaltenstil "ddmmyyyy from iso", definiert in
% preamble.tex).
Tabelle~\ref{tab:gaerung} zeigt den Verlauf der Hauptgärung anhand der
täglich protokollierten Messwerte.

\begin{table}[h]
  \centering
  \caption{Messprotokoll der Hauptgärung des Kanti Pale Ale.}
  \label{tab:gaerung}
  \footnotesize
  % Farben/Rahmen nach Vorlage tabelle-beispiel_2.html: weisse, dicke
  % Trennlinien zwischen allen Zellen (border-collapse-Optik), Kopfzeile
  % orange (bierheaderbg) mit schwarzer Schrift, Datenzeilen beige
  % (bierrowbg). Farben in preamble.tex definiert. extrarowheight sorgt
  % für mehr Abstand ober-/unterhalb des Zelltexts (auch bei mehrzeiligen
  % Bemerkung-Zellen, im Gegensatz zu \arraystretch).
  {\arrayrulecolor{white}\setlength{\arrayrulewidth}{1pt}%
  \setlength{\extrarowheight}{4pt}%
  \pgfplotstabletypeset[
    columns={Messzeitpunkt,Temperatur_Celsius,Spezifisches_Gewicht_SG,%
      Restzucker_Plato,Alkoholgehalt_Vol_Prozent,Bemerkung},
    columns/Messzeitpunkt/.style={
      ddmmyyyy from iso,
      column name={\bfseries Messzeitpunkt}, column type={|l|}},
    columns/Temperatur_Celsius/.style={
      column name={\bfseries Temp.\ (\si{\celsius})}, column type={c|}},
    columns/Spezifisches_Gewicht_SG/.style={
      column name={\bfseries Dichte (SG)}, column type={c|}},
    columns/Restzucker_Plato/.style={
      column name={\bfseries Restzucker (\si{\percent}P)}, column type={c|}},
    columns/Alkoholgehalt_Vol_Prozent/.style={
      column name={\bfseries Alkohol (\si{\percent}\,vol.)}, column type={c|}},
    columns/Bemerkung/.style={
      column name={\bfseries Bemerkung}, column type={p{3cm}|}},
    every head row/.style={before row={\hline\rowcolor{bierheaderbg}}, after row=\hline},
    every odd row/.style={before row=\rowcolor{bierrowbg}, after row=\hline},
    every even row/.style={before row=\rowcolor{bierrowbg}, after row=\hline},
    every last row/.style={after row=\hline},
  ]{daten/gaerung.csv}}%
  \par\vspace{0.3em}\footnotesize Eigene Messung, Kanti Pale Ale,
  August/Oktober 2026.
\end{table}

% Für das Diagramm werden die (bereits umformatierten) Messzeitpunkte
% zusätzlich als expliziete Tick-Label-Liste zusammengebaut: die
% Kurzform "xticklabels from table" übernimmt Zellinhalte 1:1 aus der
% Rohdatei und wendet "ddmmyyyy from iso" (ein reiner
% Typeset-/Einlese-Spaltenstil) dabei nicht an.
\pgfplotstableread[col sep=semicolon]{daten/gaerung.csv}\gaerungraw
\pgfplotstablegetrowsof{\gaerungraw}
\pgfmathtruncatemacro{\gaerungmaxrow}{\pgfplotsretval-1}
\foreach \gaerungrow in {0,...,\gaerungmaxrow}{%
  \pgfplotstablegetelem{\gaerungrow}{Messzeitpunkt}\of\gaerungraw
  \let\kswoRawDT\pgfplotsretval
  \StrLeft{\kswoRawDT}{4}[\kswoIsoY]
  \StrMid{\kswoRawDT}{6}{7}[\kswoIsoM]
  \StrMid{\kswoRawDT}{9}{10}[\kswoIsoD]
  \StrMid{\kswoRawDT}{12}{16}[\kswoIsoT]
  \edef\kswoNewDT{\kswoIsoD.\kswoIsoM.\kswoIsoY\space\kswoIsoT}
  \ifnum\gaerungrow=0
    \xdef\gaerungticklabels{\kswoNewDT}
  \else
    \xdef\gaerungticklabels{\gaerungticklabels,\kswoNewDT}
  \fi
}

Abbildung~\ref{fig:gaerverlauf} stellt denselben Datensatz grafisch dar
und zeigt, wie der Restzuckergehalt über den Gärverlauf sinkt, während
der Alkoholgehalt entsprechend ansteigt.

\begin{figure}[h]
  \centering
  \begin{tikzpicture}
    \begin{axis}[
      width=0.95\linewidth,
      height=6.5cm,
      xlabel={Messzeitpunkt},
      ylabel={Anteil (\si{\percent})},
      xtick=data,
      xticklabels/.expand once=\gaerungticklabels,
      x tick label style={rotate=60, anchor=east, font=\footnotesize},
      y tick label style={font=\footnotesize},
      label style={font=\footnotesize},
      legend style={font=\footnotesize},
      legend pos=north east,
      ymin=0, ymax=13,
      grid=major,
    ]
      \addplot table [x expr=\coordindex, y=Restzucker_Plato,
        col sep=semicolon] {daten/gaerung.csv};
      \addlegendentry{Restzucker (\si{\percent}P)}
      \addplot table [x expr=\coordindex, y=Alkoholgehalt_Vol_Prozent,
        col sep=semicolon] {daten/gaerung.csv};
      \addlegendentry{Alkohol (\si{\percent}\,vol.)}
    \end{axis}
  \end{tikzpicture}
  \caption{Verlauf von Restzucker und Alkoholgehalt während der
    Hauptgärung.}
  \label{fig:gaerverlauf}
\end{figure}

